\scanpage{19}
% archive sheet number: 18

Ceci nous amène \unsure{à|:} introduire

\[
\boxed{C_0'''}=\text{Cat. des }\deleted{\text{triples}}\inserted{\text{quadruples}}\ 
(C,Q_C,\varphi_G^{C\deleted{\unsure{G}}},h_0),
\]

\margin{\uline{NB} $n\geq 2$}
$C$ fibré projectif de rg $n$, $Q_C$ hyperquadrique dans $C$, $\varphi_G^C:G\unsure{\relbar|\longrightarrow}\operatorname{Aut}(C,Q_C)$, $h_0$ section de $C$, satisfaisant les conditions

a) les $\tau_i^C\overset{\mathrm{dfn}}=\varphi_G^C(\underline{\tau}_i)$ sont des ps. réfl., strictes si $1\leq i\leq n$, et associée à $h_0$ si $i=0$ (peut être triviale si $i=0$)
\[
\Rightarrow\ \begin{aligned}
h_i&\quad(1\leq i\leq n)\\
H_i&\quad(1\leq i\leq n)
\end{aligned}
\]
\margin{\uline{NB} $\tau_0^C$ est \unsure{str. ortho.|strict. ortho.} ssi \unsure{s'est|c'est} une réflexion stricte, i.e. ssi \deleted{$\beta_0\ne 0$} \inserted{en chaque $\mathrm{p}^\mathrm{t}$} $\forall s,\ \beta_0(s)\ne 0$ ou $\unsure{k|K}(s)$ de car. $\ne 2$}

b) Les $h_i$ lin. indép. en chaque fibre

c) $h_i\notin H_{i+1}$ en $\forall$ fibre \quad $(0\leq i\leq n-1)$

\deleted{\uline{NB} Si $Q_C$ est lisse, \inserted{et \deleted{\unsure{les}} $\tau_0$ \uline{stricte} ($\ne\mathrm{id}$ en chaque fibre)} alors \unsure{on a} \unsure{équivaut|équivat|équivalent} : $\forall i$, $H_i$ est l'orth. de \unsure{$h_i$}}

\uline{NB} Supposons les $\tau_i$ stricts orthogonaux \inserted{dans $C$ \deleted{\illegible{1}}} i.e. les $\beta_i$ inv. (C'est le cas si $Q_C$ est lisse \ldots) Alors la condition c) équivaut à

\[
\begin{aligned}
\text{c$'$)}\ (\tau_i\tau_{i+1})^2&=\mathrm{id}\quad\text{sur chaque fibre,}\\
\text{i.e.}\quad \tau_i&\ \text{ne commute pas à }\tau_{i+1}\text{ en chq fibre}
\end{aligned}
\ \Big|\ 0\leq i\leq n-1
\]

La condition b) s'énonce ainsi

b$'$) En $\forall\, s\in S$, la repr. \inserted{de $G$} induite sur $C_s$ est irréductible [suffisance claire, pour la nécessité il faut travailler un peu \ldots]
\margin{?}
